\section{Capability Boundaries and State Power}
\label{sec:boundaries}

The framework developed so far is allocative. It identifies legal subjects, permits functional classifications, maps control, and places external and internal responsibility. Allocation alone does not require every user to receive every capability or permit a government to compel every desired use. This Part identifies two boundaries that follow from the same control analysis.

First, unequal capability can be legitimate, but consequential tiering should be proportionate, transparent, reviewable, and capable of escalation. Second, the government may choose its tools and specifications, especially for national security, but procurement authority is not a general power to punish a supplier's protected refusal or conscript its technology for every lawful use. Part V has already supplied the parallel boundary for care and public service: automation can perform an existing duty, but its availability does not itself discharge one.

\subsection{Capability Stratification Without Status Stratification}

AI systems are already stratified. A public user may receive a conversational interface with limited tools. A professional may receive longer context, domain data, and auditable workflow integration. A researcher may receive weights or a sandbox. A government may receive classified hosting, specialized data, and operational tools. An enterprise may reserve higher rate limits and support for customers able to pay. These differences can reflect security, competence, scarcity, legal obligation, or commercial choice. They do not make the system a different kind of legal subject.

The first mistake is therefore to announce that unequal model access is inherently unlawful. Law regularly permits differentiated access to dangerous equipment, licensed services, secure data, and scarce resources. A medical deployment should not expose patient records and ordering tools to every public user; a government system may require a cleared environment; a research release may need controls unavailable in a consumer application. Risk-graded regulation itself relies on differentiation. The EU AI Act, for example, connects obligations to the intended use, actor, and risk of a system rather than demanding a single capability for all.\lrfootnote{See \statcite{euAIAct2024}.}

The opposite mistake is to treat every capability boundary as a private technical fact. Tiering can determine who receives an educational accommodation, who can challenge a benefits decision, which clinician sees uncertainty, or whose account is automatically banned. In those settings, the relevant provider or institution remains subject to ordinary antidiscrimination, consumer, contract, administrative, and constitutional rules. ``The model assigned the tier'' does not alter the legal subject or the nature of the decision.

Four governance conditions make stratification legible without demanding uniformity.

\begin{enumerate}[leftmargin=2.2em]
\item \term{Proportionality}: a restriction should track the material risk, resource, or legal constraint asserted. A limit designed for autonomous execution should not silently disable harmless explanation; a security rule should not become a pretext for excluding a protected class.
\item \term{Transparency}: in a consequential setting, affected persons should know that capability has been limited, the category of reason, the decisionmaker, and the route to review. Transparency does not require publication of exploit details, model weights, classified information, or another person's data.
\item \term{Review}: a person should be able to contest a mistaken tier, automated ban, or consequential capability denial before an actor with information and authority to change it. The required process depends on the underlying legal interest; a consumer feature request and termination of essential benefits are not equivalent.
\item \term{Escalation}: where broader capability can be used safely by a qualified actor or under added controls, the system should identify a path to that setting rather than treat the lowest public tier as the limit of possible service.
\end{enumerate}

These conditions derive from control rather than a new equality entitlement. NIST's framework and its generative-AI profile ask organizations to document roles, context, human oversight, credentials, feedback, impacts, and change controls; the deployment cascade identifies who can review and change a tier.\lrfootnote{See \bluecite{nistAIRMF2023}; \bluecite{nistGenAIProfile2024}. Federal management memoranda apply comparable testing, monitoring, review, acquisition, and portability expectations to executive-agency uses, without creating a general private right or rules for private platforms. See \bluecite{ombM2521}; \bluecite{ombM2522}.} Public-law requirements can add notice and hearing when government action affects a protected interest. Civil-rights law can prohibit criteria or effects within its scope. Consumer law can police a misleading representation that all users receive a capability when a hidden rule says otherwise. The minimum leaves each source of law intact and makes its object visible.

Capability and responsibility should also remain aligned. A vendor that gives a professional deployer powerful tools but withholds necessary risk information creates a different record from one that supplies evaluation, training, monitoring, and stop controls. A deployer that requests a restricted capability and then ignores its conditions cannot rely on the vendor's brand. An automated gate that no one can explain or review reveals a missing controller in the cascade. The legal question is not whether tiering exists, but whether the actor responsible for the consequential boundary can justify and correct it under the applicable rule.

\subsection{State Power: Procurement, Conditions, and Retaliation}

The state occupies two positions in AI governance. It is a regulator capable of imposing public duties, and it is an enormous purchaser capable of shaping private systems through contracts. Both powers are legitimate and bounded. The 2026 dispute between Anthropic and the United States Department of War provides the clearest current anchor.

\subsubsection{The Anthropic--Pentagon Controversy}

Anthropic developed Claude and, beginning in 2024, made versions available to federal intelligence and defense users through commercial and classified computing environments. A 2025 agreement carried a ceiling of 200 million dollars. When the government later demanded authorization for ``all lawful uses,'' Anthropic substantially narrowed its usage restrictions but preserved two contractual red lines: lethal autonomous warfare and mass surveillance of Americans. The deployed models were static; Anthropic could not see government prompts, technically block a particular use, or reach the systems through a back door. The parties reached no agreement by the government's February 27 deadline, although negotiation and collaboration continued afterward. According to the district court's record, three challenged actions followed. A February 27 Presidential Directive ordered government-wide cessation with a six-month transition. Later that day, the Secretary of Defense directed commencement of a supply-chain designation process while immediately announcing a broader prohibition reaching military contractors and suppliers. A March 3 determination, noticed March 4, then formally invoked 10 U.S.C. \S~3252 for national-security-system procurement.\lrfootnote{See \casecite{anthropicFinal2026}. The record does not establish that Claude was deployed in Project Maven; this Article accordingly uses the case caption and the Anthropic--Pentagon controversy rather than treating ``Claude--Maven'' as an adjudicated fact.}

In March 2026, Judge Rita F. Lin preliminarily enjoined the directives and designation. The order emphasized a boundary that survives every procedural stage: it did not require the government to use Anthropic's products and did not prevent a lawful transition to another provider.\lrfootnote{See \casecite{anthropicPI2026}.} On August 27, 2026, the district court granted summary judgment to Anthropic on its First Amendment retaliation and Fifth Amendment stigma-plus due-process claims and on core Administrative Procedure Act claims, entered final judgment, and closed the case. It held that 10 U.S.C. section 3252 addressed adversarial sabotage or subversion, not punishment of a domestic supplier for its public safety position; the designation was beyond authority, procedurally deficient, arbitrary, and pretextual. It also found no statutory authority for the Secretary's secondary boycott of firms doing business with the military. The court vacated the designation and barred enforcement of the retaliatory directives while preserving ordinary lawful procurement choice.\lrfootnote{See \casecite{anthropicFinal2026}; \casecite{anthropicRelief2026}. This was an appealable district-court judgment, not a final appellate disposition; it remained subject to appellate review as of this Article's date. A separate D.C. Circuit proceeding concerning 41 U.S.C. sections 4713 and 1327 presents distinct statutory and procedural questions; its April 2026 emergency-stay ruling did not decide the merits. See \casecite{anthropicDCStay2026}.}

The controversy is important because the strongest argument for the government is substantial. Military purchasers must control mission requirements, reliability, availability, classified integration, and lawful use. A vendor's unilateral restriction can create operational risk, and the government need not design a mission around one company's policy. National-security procurement also involves predictive judgments that courts are institutionally cautious to second-guess. Nothing in the jurisprudential minimum gives an AI provider a right to a federal contract.

That argument supports selection, specification, and transition. It does not establish a power to impose market-wide penalties because a company criticizes the government's position or refuses to alter its general product. The final opinion distinguished the government as market participant administering a contract from the government as sovereign imposing a public stigma, government-wide ban, and third-party boycott. Government contractors retain First Amendment protection against retaliation for protected expression, subject to the government's interests as contractor and purchaser. \emph{Board of County Commissioners v. Umbehr} applied that principle to termination or nonrenewal of an independent contractor, and \emph{O'Hare Truck Service, Inc.\ v. City of Northlake} rejected removal from a rotation list based on political association.\lrfootnote{See \casecite{umbehr1996}; \casecite{ohare1996}.} The unconstitutional-conditions doctrine more generally prevents government from denying a benefit on a basis that infringes a protected interest even where the claimant had no entitlement to the benefit itself.\lrfootnote{See \casecite{perrySindermann1972}; \casecite{aosi2013}. This Article uses unconstitutional conditions as an analytic bridge; the district court's final holdings sounded directly in retaliation, due process, and the APA.}

The procurement boundary can be stated as a sequence. It is this Article's synthesis of purchaser discretion, contractor-retaliation doctrine, and unconstitutional-conditions law, not a four-part test announced by the district court. Because the sequence rests on \emph{Umbehr}, \emph{O'Hare}, \emph{Perry}, and the distinction between purchasing and sovereign sanction, it does not depend on the August 27 judgment remaining undisturbed in every respect:

\begin{enumerate}[leftmargin=2.2em]
\item The government may define the mission, security conditions, performance requirements, and lawful uses for which it seeks a system.
\item A supplier may decide which product it offers and may express a view about uses, subject to contracts and law that validly bind it.
\item The government may select another supplier, negotiate different terms, build its own system, or take a tailored action authorized by statute for a genuine security or performance risk.
\item It may not use procurement as a pretext to punish protected criticism, compel adoption of a viewpoint outside the funded program, or impose sweeping third-party disabilities unsupported by the invoked authority.
\end{enumerate}

This is not a special privilege for safety-branded companies. A vendor cannot convert every commercial limit into protected speech, abandon a promised service without contractual consequence, or invoke ethics to conceal a security defect. The court should identify the expressive activity, adverse action, causal evidence, statutory authority, and government's purchaser interests under ordinary doctrine. The \emph{Anthropic} record mattered because the stated sanctions extended far beyond choosing a different model and the court found the asserted security rationale unsupported.

The subject-status invariant clarifies the dispute. Claude was not the constitutional claimant; Anthropic was. The government was not bargaining with an artificial rights holder; it was setting terms with a corporation that designed and maintained a system. The contested constraints embodied human and organizational choices. Protecting or rejecting those choices required no ruling on whether Claude's generated output itself possessed First Amendment rights.

The same principle constrains providers. A company cannot invoke its own model constitution to override positive law, evade a valid contract, or deny public remedies. Internal policy is evidence of design and commitment; external law judges it. The framework is therefore reciprocal. It prevents the firm from transferring responsibility downward to the model, and it prevents the state from treating control over procurement as control over every lawful private use or viewpoint.

Taken together, capability stratification and bounded state power complete the architecture's reciprocal account of control. The framework allocates authority to those who actually hold it and preserves the legal limits that accompany that authority.
